% =============================================================================
% Chapter 1  Introduction  --  ~3.5 pp  --  rubric: Introduction & Literature
% Review /20.   Plan: ThesisPlanning.md §5.1.
%
% PURPOSE     Establish that quantum DE solvers matter and that the physics-
%             informed Hamiltonian route has a specific, *located* cost problem;
%             then state what this thesis contributes, and what it does not.
% CRITERION   "Comprehensive, superior understanding of the historical background
%             and state of the art, with excellent connection of the current field
%             to the project" (>= 18/20). The introduction converges; it does not
%             survey.
% SOURCES     Planning.md §0; ThesisPlanning.md §3.1 (contributions), §3.2 (spine
%             sentence), §3.3 (vocabulary).
% WRITE WHEN  Last (wave W6), so it promises only what was delivered.
% =============================================================================

\subsection{Differential Equations as a Target for Quantum Computation}
\label{sec:intro-target}
% ~1 pp
% LAND, in order:
%   1. Applications (fluids, finance, fields), with citations to the domains.
%   2. Why quantum algorithms are proposed for them.
%   3. The resource that matters: queries x gates per query, and qubits.
%   4. The algorithm families in one or two sentences each; the detail and the
%      judgement are Chapter 2's job (sec:background).

\subsection{Physics-Informed Hamiltonians and Where Their Cost Lies}
\label{sec:intro-cost}
% ~1 pp
% LAND, in order:
%   1. Wu et al.'s idea in three sentences, attributed at FIRST mention
%      (\cite{wu2025pihm}), not at first critique.
%   2. The question this thesis asks: is the method's cost in its ENCODER or in
%      its HAMILTONIAN?
%   3. One quantitative sentence before the end of page 2. At n = 8 (the paper's
%      Fig. 3b) the filter degree is 8.2e23 as published, 1.2e10 with the best
%      exact encoding, and 5.6e3 in descriptor form (Planning.md App. A-6).
%      Quote with \res once pihm produces them (\prov{...} until then); check
%      which filter the final numbers use (A-6 quotes the minimax edge filter).

\subsection{Contributions and Scope}
\label{sec:contributions}
% ~1 pp
% LAND, in order:
%   1. The numbered contributions (ThesisPlanning.md §3.1): the headline C, then
%      S1-S4. Each is one sentence of WHAT and one clause of EVIDENCE (a chapter
%      pointer). C always carries its cost (the block register); S2 is phrased
%      as serving the comparison ("isolates ..."), never as a rival solver.
%        C   the descriptor reformulation                      -> Ch. 5, §8.3-8.4
%        S1  a verified reproduction and audit of Wu et al.     -> §4.1, §8.1, App. K
%        S2  an exact structured encoding of the standard form  -> Ch. 4, §8.3
%        S3  a regime-independent verified pipeline             -> Ch. 7, §8.2
%        S4  nonlinearity: preparation versus projection        -> Ch. 6, §8.5-8.6
%   2. The spine sentence (ThesisPlanning.md §3.2).
%   3. The scope statement, once, plainly: every result is a classical
%      computation -- circuit simulation (T1) at small n, exact emulation (T2)
%      beyond it, resource estimates (T3) beyond that. Nothing ran on hardware.
%      Readout is classical, decoded from the prepared state's field blocks; the
%      paper's interferometric protocol is built and verified only at small n
%      (R1, R2a, R5a), with our correction for p_G < 1, and is further work.
% REVISE (T-24): the line above replaces "Readout is the paper's protocol,
%   reproduced in simulation", which overstates what is reproduced and
%   understates that the thesis's own readout is classical decoding, not the
%   paper's interferometric protocol (CLAUDE.md Standing Facts, Readout).
%   4. What is NOT claimed: a poly-log end-to-end cost; preparation for
%      degenerate formulations; a new readout protocol.

\subsection{Outline}
\label{sec:outline}
% ~0.5 pp
% LAND: one item per chapter naming its job, headings quoted verbatim. Mark
%   Chapter 3 as reviewed prior art and Chapters 4-7 as this thesis's model.

% CHECKLIST (marker)
%   1. Can a non-specialist physicist state the problem after two pages?
%   2. Is there a number on or before page 2?
%   3. Is every contribution numbered and falsifiable, with its evidence named?
%   4. Is the scope statement explicit, including the tier vocabulary and
%      "no hardware"?
%   5. Is Wu et al. attributed at first mention, not at first critique?
%   6. Does the outline match the headings verbatim?
% TRAPS  Leading with the descriptor's advantage before the reader knows the
%        baseline; claiming a speed-up (Style Trap 9); listing S2 as a competing
%        method.
